% cap01/1.3_hipotesis.tex — Hipótesis de trabajo.
\section{Hipótesis}
\label{section:hipotesis}
Conforme a Hernández-Sampieri y Mendoza (2018), en los estudios con alcance descriptivo la formulación de hipótesis se considera pertinente y necesaria cuando la investigación pronostica el comportamiento de un hecho o variable de diseño. En estricta coherencia con ello, y sustentado en el diseño teórico-matemático y topológico del sistema de software, se enuncia la siguiente hipótesis de trabajo de naturaleza contrastable:

\begin{quote}
\itshape La arquitectura de la plataforma Democra, mediante la aplicación implacable de políticas de aislamiento de seguridad a nivel de fila (\textit{Row Level Security}, RLS) y un control de acceso basado en roles (RBAC) diseñado bajo el paradigma de negación por defecto, garantiza matemáticamente el aislamiento absoluto de los datos entre inquilinos (logrando cero accesos \textit{cross-tenant} no autorizados). Adicionalmente, se plantea que esta robustez \textit{backend} no degrada la experiencia de usuario, y que la plataforma alcanzará un nivel de funcionamiento verificable de sus módulos críticos, validado mediante pruebas automatizadas y métricas de cobertura de código que garanticen la correcta operación del sistema.
\end{quote}

La contrastación empírica de esta hipótesis metodológica se operacionaliza exhaustivamente en el \autoref{chapter:metodo}, y los hallazgos derivados de las métricas e instrumentación analítica se documentan a detalle en el \autoref{chapter:resultados_discusion}.
